\begin{table}[tp]
\centering
\caption{Finite-sum smooth convex optimization: \href{https://github.com/k3nfalt/longtails/blob/main/LongTails/Setups/FiniteSumSmoothAndConvex.lean}{FiniteSumSmoothAndConvex}.}\label{tab:FiniteSumSmoothAndConvex}
\parbox{\linewidth}{\footnotesize $F=n^{-1}\sum_{i=1}^n f_i$ on a real Hilbert space $E$ equipped with a measurable structure; $n\ge1$, each $f_i$ is convex with $L$-Lipschitz gradient, $L>0$. A minimizer $x^\star$ satisfies $\|x_0-x^\star\|\le R$, $R\ge0$. One oracle call returns an exact component gradient. Target: measurable output $\widehat x$ with integrable gap and $\int [F(\widehat x(\omega))-F(x^\star)]\,dP(\omega)\le\varepsilon$, $\varepsilon>0$, where $P$ is the algorithm's sampling law. The setup imposes no algorithm restriction.}\par\smallskip
\begin{threeparttable}
\footnotesize
\setlength{\tabcolsep}{4pt}\renewcommand{\arraystretch}{1.25}
\begin{tabularx}{\linewidth}{@{}>{\raggedright\arraybackslash}p{0.20\linewidth}>{\raggedright\arraybackslash}p{0.32\linewidth}>{\raggedright\arraybackslash}p{0.25\linewidth}>{\raggedright\arraybackslash}X@{}}
\toprule
\textbf{Result} & \textbf{Oracle complexity} & \textbf{Novelty} & \textbf{Limitations} \\
\midrule
\textbf{(1)} \href{https://github.com/k3nfalt/longtails/blob/main/LongTails/Theorems/AcceleratedFiniteSumConvex.lean}{Accelerated\allowbreak Finite\allowbreak Sum\allowbreak Convex} \textbf{(Ours)} \href{https://github.com/k3nfalt/longtails/blob/main/scripts/reviews/AcceleratedFiniteSumConvex.md}{(review)}
& $\mathcal{O}\!\left(n[1+\log\log(n+2)]+\sqrt{nLR^2/\varepsilon}\right)$ & & \\
\midrule
\textbf{(2)} \href{https://github.com/k3nfalt/longtails/blob/main/LongTails/Theorems/SvrgConvex.lean}{Svrg\allowbreak Convex} \textbf{(Ours)} \href{https://github.com/k3nfalt/longtails/blob/main/scripts/reviews/SvrgConvex.md}{(review)}
& $\mathcal{O}\!\left(1+n\log(1+LR^2/\varepsilon)+LR^2/\varepsilon\right)$ & & \\
\midrule[1.2pt]
\textbf{(3)} VRADA \citep[Thm.~1, (9)--(11)]{finiteSumSong2021v4}. $E=\mathbb R^d$; $l=0$, $\sigma=0$, $m=n$.
& $\mathcal{O}\!\left(n[1+\log\log(n+2)]+\sqrt{nLR^2/\varepsilon}\right)$
& \textbf{(1)}: Matches known rate; Lean formalization on Hilbert spaces. \textbf{(2)}: Smaller startup when $1\ll T\ll\log n$.
& \textbf{(2)}: $T$ instead of $\sqrt{nT}$ for high accuracy. \textbf{(1)}, \textbf{(2)} do not cover the additional composite objective supported by \textbf{(3)}. \\
\midrule
\textbf{(4)} ANITA \citep[Thm.~1, Cor.~1]{finiteSumLi2022v3}. $E=\mathbb R^d$.
& $\displaystyle\mathcal{O}\!\left(\begin{aligned}&n\!\left[1+\min\!\left\{\log_+\!\frac{LR^2}{\varepsilon\sqrt n},\log\sqrt n\right\}\right]\\[-2pt]&{}+\sqrt{nLR^2/\varepsilon}\end{aligned}\right)$\par Expected call count.
& \textbf{(1)}: Intermediate-accuracy savings for $\sqrt n\log n\lesssim T\lesssim n\log^2 n$; at $T=n$, $\mathcal{O}(n\log\log n)$ versus \textbf{(4)}'s $\mathcal{O}(n\log n)$ for the shared expected-gap target.
& \textbf{(1)}: At $T\le\sqrt n$, $\mathcal{O}(n\log\log n)$ versus \textbf{(4)}'s $\mathcal{O}(n)$. Epoch average instead of last iterate; $\varepsilon$ needed for stopping, while \textbf{(4)}'s updates are accuracy-independent. \\
\midrule
\textbf{(5)} SVRG++ \citep[Thm.~4.1]{finiteSumAllenZhuYuan2016}. $E=\mathbb R^d$, $\Psi=0$, $m_0=1$.
& $\mathcal{O}\!\left(n[1+\log(1+LR^2/\varepsilon)]+LR^2/\varepsilon\right)$
& \textbf{(1)}: At $T=n$, $\mathcal{O}(n\log\log n)$ versus \textbf{(5)}'s $\mathcal{O}(n\log n)$; accuracy term $\sqrt{nT}$ instead of $T$ when $T>n$. \textbf{(2)}: Matches known rate; dimension-free Lean formalization.
& \textbf{(1)}, \textbf{(2)} omit the composite term supported by \textbf{(5)}. \\
\midrule
\textbf{(6)} SVRG++ \citep[Thm.~4.1]{finiteSumAllenZhuYuan2016}. $E=\mathbb R^d$, $\Psi=0$; $m_0=\lceil LR^2/\Delta\rceil$, $0<\varepsilon<\Delta$.
& $\mathcal{O}\!\left(n[1+\log(\Delta/\varepsilon)]+LR^2/\varepsilon\right)$
& \textbf{(2)}: Matches known rate when $\Delta\asymp LR^2$.
& \textbf{(2)}: Uses $\log(1+LR^2/\varepsilon)$ instead of $1+\log(\Delta/\varepsilon)$; can lose when $\Delta\ll LR^2$. No composite term. \\
\midrule
\textbf{(7)} Katyusha + AdaptReg \citep[Cor.~3.5]{AllenZhu2018KatyushaComposite}. $E=\mathbb R^d$, $\psi=0$; $0<\varepsilon<\Delta$.
& $\mathcal{O}\!\left(n[1+\log(\Delta/\varepsilon)]+\sqrt{nLR^2/\varepsilon}\right)$
& \textbf{(2)}: No established improvement.
& \textbf{(2)}: $T$ instead of $\sqrt{nT}$; up to $\sqrt{T/n}$ higher total cost for $T\gg n$. \textbf{(1)}, \textbf{(2)} omit the composite term supported by \textbf{(7)}. \\
\bottomrule
\end{tabularx}
\begin{tablenotes}[flushleft]\footnotesize
\item[] \textbf{Abbr:} $n$: component count; $L$: component smoothness; $R$: initial-distance bound; $\varepsilon$: expected-gap tolerance; $T=LR^2/\varepsilon$; $\Delta=F(x_0)-F(x^\star)\le LR^2/2$; $\log_+u=\max\{0,\log u\}$ (zero at $u=0$). Costs count component gradients, including full-gradient refreshes; \textbf{(4)} counts calls in expectation. $l,\Psi,\psi$: optional convex composite terms; $\sigma$: strong-convexity parameter; $m,m_0$: epoch-length parameters.
\end{tablenotes}
\end{threeparttable}
\end{table}
\clearpage
